\chapter{Videos, música y fotos}\label{cap:10-multimedia}

\section{YouTube ligero}\label{sec:youtube-ligero}

Ver YouTube en el navegador gasta mucho procesador: en equipos viejos el
video se traba aunque la conexión sea buena. \textbf{YouTube ligero}
busca los videos en una ventana sencilla (\emph{Pipe Viewer}) y los
reproduce con \emph{mpv}, un reproductor muy liviano, en 480p, el formato
que menos procesador gasta \cite{mpv,ytdlp}. No necesita navegador ni
cuenta.

\begin{pasos}
  \item Abre \menu{Aplicaciones > Multimedia > YouTube}.
  \item Escribe lo que buscas arriba y pulsa \keys{Enter}.
  \item Haz doble clic en un video. Se abre en una ventana de mpv.
\end{pasos}

\captura{youtube-ligero}{YouTube ligero: la búsqueda en Pipe Viewer}

Mientras se reproduce: \keys{Espacio} pausa, \keys{\arrowkeyleft} y
\keys{\arrowkeyright} retroceden y adelantan 5 segundos, \keys{F} pone
pantalla completa, \keys{9} y \keys{0} bajan y suben el volumen, \keys{Q}
cierra.

\begin{description}
  \item[Si el video se traba] (procesadores de un solo núcleo): en Pipe
    Viewer, abre el menú de preferencias y elige resolución \emph{360}. Se
    guarda al cerrar la ventana.
  \item[Si deja de funcionar] YouTube cambia seguido y a veces rompe los
    programas que lo leen. Instala las actualizaciones
    (sección~\ref{sec:actualizar}); si no alcanza, abre una terminal y
    escribe:
\begin{terminal}
doas apk upgrade yt-dlp yt-dlp-ejs pipe-viewer
\end{terminal}
\end{description}

\section{Videos y música}

Lázaro abre los videos y los audios con \textbf{mpv}: haz doble clic en el
archivo. Reproduce casi cualquier formato, incluidos los audios de
WhatsApp. Las teclas son las mismas de YouTube ligero.

Si al instalar marcaste \textbf{VLC}, tienes un reproductor más completo,
con menús, listas de reproducción y subtítulos \cite{vlc}. Para abrir un
archivo con VLC: clic derecho, \menu{Abrir con > VLC}. Para que sea el
predeterminado, en esa misma ventana marca la casilla para usarlo siempre.

\begin{consejo}
¿Una película con subtítulos en un archivo aparte (\texttt{.srt})? Ponle el
mismo nombre que al video y déjalo en la misma carpeta: mpv y VLC lo
cargan solos.
\end{consejo}

\subsection{El sonido}
El icono del altavoz en la bandeja sube y baja el volumen; con clic
derecho (o \menu{Configuración del sonido}) se elige la salida: los
parlantes, los audífonos, la pantalla por HDMI o unos audífonos
Bluetooth. Si no se escucha nada, mira la sección~\ref{sec:prob-sonido}.

\section{Fotos}

Haz doble clic en una foto para verla. Con las flechas pasas a la
siguiente. En XFCE se abre con \emph{Ristretto}; en los demás escritorios,
con su visor propio.

\begin{description}
  \item[Pasar las fotos del celular] Conecta el celular por cable. En
    Android, desliza la barra de arriba, toca la notificación del USB y
    elige \emph{Transferencia de archivos}. El celular aparece en el
    gestor de archivos como un disco más (capítulo~\ref{cap:11-archivos}).
  \item[Retocar] Para recortar, girar o mejorar el color, la tienda
    ofrece programas como \emph{GIMP} (completo) o \emph{Pinta} (sencillo).
\end{description}

\section{Capturas de pantalla}

En XFCE, la tecla \keys{ImprPant} abre \emph{Captura de pantalla}: elige si
quieres toda la pantalla, una ventana o un área, y dónde guardarla. En
GNOME y KDE Plasma esa misma tecla abre su herramienta propia.

\section{Cámara web}

Para probar la cámara, abre una videollamada de prueba en el navegador
(por ejemplo, \url{https://meet.jit.si}). Si la página pide permiso para
usar la cámara, pulsa \menu{Permitir}.

\fuentes
